\chapter*{Introduction}
\markleft{Introduction}
\addcontentsline{toc}{chapter}{Introduction}

L'institut AV-TEST recense chaque jour plus de 450\,000 nouveaux logiciels
malveillants \cite{avtest}. Le chiffre impressionne moins qu'il n'y paraît. Ces
échantillons ne sont pas 450\,000 menaces distinctes : ce sont, pour l'essentiel,
des variantes produites par des moteurs de mutation dont la logique demeure finie.
Un analyste qui reconstitue cette logique neutralise la famille entière. La
détection par signature tient sur ce pari.

Le rapport \emph{AI Threat Tracker} publié le 5 novembre 2025 par le Google Threat
Intelligence Group documente les premiers programmes qui ne le tiennent plus
\cite{gtigAI}. GTIG y décrit des \emph{malwares} qui interrogent un modèle de
langage pendant leur exécution, pour générer des scripts, obscurcir leur propre
code ou produire des fonctions malveillantes à la demande. PROMPTSTEAL, employé
contre des cibles ukrainiennes, interroge un modèle hébergé sur Hugging Face pour
générer ses commandes Windows au lieu de les coder en dur. PROMPTFLUX sollicite
l'API Gemini pour se réécrire « \emph{just-in-time} ». PROMPTLOCK, rançongiciel,
génère à l'exécution les scripts Lua qui parcourent le système de fichiers,
sélectionnent les données et les chiffrent. L'intelligence artificielle générative
n'intervient plus seulement en amont, dans la rédaction de courriels de
\emph{phishing} ou l'écriture assistée de charges utiles : elle est appelée par le
programme lui-même, sur la machine de la victime.

Jusqu'ici, le polymorphisme et le métamorphisme restaient déterministes. Le code
mutait, mais selon des règles finies, qu'un analyste pouvait reconstituer et
généraliser. C'est ce postulat que la génération à l'exécution abolit. Il n'y a
plus de moteur de mutation à rétro-concevoir, seulement une invite et un modèle
dont les sorties varient d'une itération à l'autre.

Encore faut-il le vérifier ce comportement plutôt que lire ces rapports. Ce mémoire reprend le
dernier de ces trois échantillons et l'analyse statiquement, puis dynamiquement,
avant de convertir les observations en règles YARA, en règles Sigma et en
\emph{plugin} Volatility. Ce parcours, du binaire à la règle, est celui que suit un
analyste en SOC lorsqu'un échantillon inconnu lui parvient.

Il conduit à la question qui structure ce travail. Nos mécanismes de détection
supposent un adversaire dont le comportement demeure \emph{in fine} modélisable.
Que valent-ils face à un code qui change à chaque itération et s'optimise contre
les défenses qu'il rencontre ?